% from ega4-17-fr.tex
\oldpage[IV-4]{57}
  est surjective (resp. injective, resp. bijective); par composition, on en conclut qu'il en est de même de (17.1.1.1).
  \item[(iii)] On notera que les propriétés du morphisme $f$ définies dans (17.1.1) sont des propriétés du \emph{foncteur représentable} $(0_{\mathrm{III}}, 8.1.8)$
  \[
    Y'\longmapsto\Hom_Y(Y',X)
  \]
  de la catégorie des $Y$-préschémas dans la catégorie des ensembles; elles gardent un sens pour \emph{n'importe quel} foncteur contravariant ayant mêmes catégories de départ et d'arrivée, représentable ou non.
  \item[(iv)] Supposons que le morphisme $f$ soit formellement non ramifié (resp. formellement étale); considérons un $Y$-préschéma \emph{quelconque} $Z$ et un sous-préschéma fermé $Z_0$ de $Z$ défini par un Idéal \emph{localement nilpotent} $\sh{J}$ de $\sh{O}_Z$. Alors l'application
  \[
  \label{IV.17.1.2.1}
    \Hom_Y(Z,X)\longrightarrow\Hom_Y(Z_0,X)
    \tag{17.1.2.1}
  \]
  déduite de l'injection canonique $Z_0\to Z$, est encore injective (resp. bijective). En effet, soit $(U_\alpha)$ un recouvrement ouvert affine de $Z$ tel que les Idéaux $\sh{J}|U_\alpha$ soient nilpotents, et pour tout $\alpha$, soit $U^0_\alpha$ l'image réciproque de $U_\alpha$ dans $Z_0$, qui est le sous-schéma fermé de $U_\alpha$ défini par $\sh{J}|U_\alpha$. Soit $f_0:Z_0\to X$ un $Y$-morphisme; par hypothèse, pour tout $\alpha$, il y a au plus un (resp. un et un seul) $Y$-morphisme $f_\alpha:U_\alpha\to X$ dont la restriction à $U^0_\alpha$ soit égale à $f_0|U^0_\alpha$. On en conclut aussitôt que si $f_\alpha$ et $f_\beta$ sont définis, alors, pour tout ouvert affine $V\subset U_\alpha\cap U_\beta$, on a $f_\alpha|V=f_\beta|V$, puisque les restrictions de ces morphismes à l'image réciproque $V_0$ de $V$ dans $Z_0$ coïncident. Il y a donc au plus un (resp. un et un seul) $Y$-morphisme $f:Z\to X$ dont la restriction à $Z_0$ coïncide avec $f_0$.
\end{enumerate}
\end{remarks}

\begin{proposition}[17.1.3]
\label{IV.17.1.3}
\begin{enumerate}
  \item[(i)] Un monomorphisme de préschémas est formellement non ramifié; une immersion ouverte est formellement étale.
  \item[(ii)] Le composé de deux morphismes formellement lisses (resp. formellement non ramifiés, resp. formellement étales) est formellement lisse (resp. formellement non ramifié, resp. formellement étale).
  \item[(iii)] Si $f:X\to Y$ est un $S$-morphisme formellement lisse (resp. formellement non ramifié, resp. formellement étale), il en est de même de $f_{(S')}:X_{(S')}\to Y_{(S')}$ pour toute extension $S'\to S$ du préschéma de base.
  \item[(iv)] Si $f:X\to X'$ et $g:Y\to Y'$ sont deux $S$-morphismes formellement lisses (resp. formellement non ramifiés, resp. formellement étales), il en est de même de $f\times_S g:X\times_S Y\to X'\times_S Y'$.
  \item[(v)] Soient $f:X\to Y$, $g:Y\to Z$ deux morphismes; si $g\circ f$ est formellement non ramifié, il en est de même de $f$.
  \item[(vi)] Si $f:X\to Y$ est un morphisme formellement non ramifié, il en est de même de $f_{\mathrm{red}}:X_{\mathrm{red}}\to Y_{\mathrm{red}}$.
\end{enumerate}
\end{proposition}

En vertu de (I, 5.5.12), il suffit de prouver (i), (ii) et (iii). Les deux assertions de (i) sont triviales. Pour prouver (ii), considérons deux morphismes $f:X\to Y$, $g:Y\to Z$, un schéma affine $Z'$, un sous-schéma fermé $Z'_0$ de $Z'$ défini par un Idéal nilpotent et un morphisme $Z'\to Z$. Supposons $f$ et $g$ formellement lisses, et considérons un
